\documentclass[10pt,a4paper]{amsart}
\usepackage[margin=19mm]{geometry}
\usepackage{amsmath,amssymb,mathtools}
\usepackage[scr=rsfs,cal=euler]{mathalfa}
\usepackage{xfrac}
\usepackage[unicode,hidelinks,bookmarks=false]{hyperref}
\newcommand{\ie}{i.e.\ }
\newcommand{\cf}{cf.\ }
\newcommand{\resp}{respectively\ }
\newcommand{\Subsection}{\subsection}
\providecommand{\nobreakdash}{\nobreak\hbox{-}}
\setlength{\emergencystretch}{2em}
\multlinegap0pt
\usepackage{amssymb}
\usepackage{mathtools}
\usepackage{tcolorbox}
\usepackage{xcolor}
\newtheorem{assumption}{Assumption}
\newtheorem{lemma}{Lemma}
\newtheorem{theorem}{Theorem}
\newcommand{\dd}{\mathrm{d}}
\title{State-space models through the lens of ensemble control}
\author{Ye Feng, Jianfeng Lu}
\date{Source: arXiv:2603.13587v2 (CC BY 4.0)}
\begin{document}
\maketitle

\noindent\emph{Source and licence.} State-space models through the lens of ensemble control. Version arXiv:2603.13587v2 (CC BY 4.0), distributed by arXiv under the Creative Commons Attribution 4.0 International licence, http://creativecommons.org/licenses/by/4.0/, as recorded in the arXiv licence metadata for this exact version and shown on the abstract page https://arxiv.org/abs/2603.13587v2. Full licence notice: \texttt{licenses/arxiv-2603.13587-CC-BY-4.0.txt}.\par\smallskip

\noindent\textit{Editorial scope.} Counted unit: the article's theorem thm:ensemble-relative-smoothness (source0.tex lines 1340--1372). The article's complete original proof of it is delivered with it (source0.tex lines 1374--1456, one \texttt{proof} environment of 83 lines). No mathematics is written, repaired or re-derived: the statement and the proof are the article's own lines, in the article's own notation, and no retained line is substituted or re-flowed.  Retained uncounted as the source-local context the proof needs: lines 536--589; lines 726--738; lines 774--790; lines 971--984; lines 1184--1192.  Nothing was omitted from any retained span, so the package carries no omission record.  No retained line contains a citation, so the wrapper carries no bibliography and no bibliography entry.  No retained line contains a figure environment, an image-inclusion command or drawing code, so no image is delivered and the package declares no asset dependency.  Editorial formatting only, in addition: an AMS \texttt{amsart} preamble that restates the article's own declarations and macros used by the retained lines, namely the environments \texttt{assumption}, \texttt{lemma}, \texttt{theorem} on separate counters, together with the standard packages those lines need.  The retained environments print in this document as Assumption 1, Lemma 1, Theorem 1 in order of appearance, whereas the article prints the same items with the numbering of its own full document.  The complete native source of the article as distributed in the arXiv e-print for this exact version is bundled unchanged as \texttt{sources/196341/source0.tex}.
\begin{assumption}[Regularity of the general ensemble problem]
\label{ass:general-first-variation}
The following conditions hold.
\begin{enumerate}
    \item $U$ is a nonempty compact convex subset of $\mathbb R^m$ with nonempty interior.
    \item $x_0:\Omega\to\mathbb R^n$ is measurable and uniformly bounded;
    write $\|x_0\|_\infty:=\sup_{\omega\in\Omega}|x_0(\omega)|$.
    \item For every $(x,\theta)$${}\in\mathbb R^n\times U$, the maps
    \[
        (t,\omega)\mapsto b(t,x,\theta,\omega),\qquad
        (t,\omega)\mapsto f(t,x,\theta,\omega)
    \]
    are measurable. There is a constant $M_b>0$ such that
    \[
        |b(t,x,\theta,\omega)|\le M_b(1+|x|),
        \qquad x\in\mathbb R^n,\quad\theta\in U.
    \]
    \item There is an open neighborhood $\mathcal O_U\subset\mathbb R^m$
    of $U$ such that, for each $(t,\omega)$, the maps
    $(x,\theta)\mapsto b(t,x,\theta,\omega)$ and
    $(x,\theta)\mapsto f(t,x,\theta,\omega)$ are twice continuously
    differentiable on $\mathbb R^n\times\mathcal O_U$.
    For every $R>0$ there is $M_R>0$, independent of $(t,\omega)$, such that
    for $|x|\le R$ and $\theta\in U$,
    \[
        \begin{aligned}
        &\|D_x b(t,x,\theta,\omega)\|_{\mathrm{op}}
        +\|D_\theta b(t,x,\theta,\omega)\|_{\mathrm{op}}
        \\
        &\qquad+\|D^2_{(x,\theta)}b(t,x,\theta,\omega)\|_{\mathrm{op}}\le M_R,
        \end{aligned}
    \]
    and
    \[
        \begin{aligned}
        &|f(t,x,\theta,\omega)|+|\nabla_xf(t,x,\theta,\omega)|
        \\
        &\qquad+|\nabla_\theta f(t,x,\theta,\omega)|
        +\|D^2_{(x,\theta)}f(t,x,\theta,\omega)\|_{\mathrm{op}}\le M_R.
        \end{aligned}
    \]
    \item For each $\omega$, $g(\cdot,\omega)$ is twice continuously
    differentiable. For every $R>0$ there is $M_R>0$, independent of
    $\omega$, such that
    \[
        |g(x,\omega)|+|\nabla_xg(x,\omega)|
        +\|\nabla_x^2g(x,\omega)\|_{\mathrm{op}}\le M_R,\qquad |x|\le R.
    \]
    For each $x$, the functions $g(x,\cdot)$, $\nabla_xg(x,\cdot)$,
    and $\nabla_x^2g(x,\cdot)$ are measurable.
    \item $h$ is convex and admits a twice continuously differentiable
    extension to an open neighborhood of $U$.
\end{enumerate}
\end{assumption}

\begin{equation}
\label{eq:general-Hamiltonian-u-derivative}
\nabla_\theta H_t^\tau(\bar\xi,\bar\pi,\theta)
=
\mathbb E_\omega\!\left[
D_\theta b(t,\bar\xi(\omega),\theta,\omega)^\top
\bar\pi(\omega)
-
\nabla_\theta f(t,\bar\xi(\omega),\theta,\omega)
\right]
-
\tau\nabla h(\theta).
\end{equation}

\begin{lemma}[First variation of the ensemble objective]
\label{lem:first-variation-general}
Suppose Assumption~\ref{ass:general-first-variation} holds. For all
$u,v\in\mathcal U[0,T]$, the one-sided feasible directional derivative
\[
dJ^\tau[u](v-u):=\lim_{\varepsilon\downarrow0}
\frac{J^\tau[u+\varepsilon(v-u)]-J^\tau[u]}{\varepsilon}
\]
exists and satisfies
\begin{equation}
\label{eq:general-first-variation}
dJ^\tau[u](v-u)
=-\int_0^T\left\langle
\nabla_\theta H_t^\tau(x^u(t,\cdot),p^u(t,\cdot),u(t)),
v(t)-u(t)\right\rangle\,\dd t.
\end{equation}
\end{lemma}

\begin{assumption}[Uniform convexity of the mirror map]
\label{ass:mirror-map}
There exists $\sigma_h>0$ such that
\begin{equation}
\label{eq:uniform-convexity-h}
D_h(\vartheta\mid\theta)\ge\frac{\sigma_h}{2}|\vartheta-\theta|^2,
\qquad\theta,\vartheta\in U.
\end{equation}
Consequently, for admissible $u,v$,
\begin{equation}
\label{eq:integrated-uniform-convexity}
\mathcal D_h[v\mid u]\ge\frac{\sigma_h}{2}\|v-u\|_{L^2}^2.
\end{equation}
\end{assumption}

\begin{equation}
\label{eq:ensemble-convergence-bounds}
B_\theta:=\sup\|D_\theta b\|_{\mathrm{op}},
\qquad
L_{ij}^{b}:=\sup\|D_iD_jb\|_{\mathrm{op}},
\qquad
L_{ij}^{f}:=\sup\|D_iD_jf\|_{\mathrm{op}},
\quad i,j\in\{x,\theta\}.
\end{equation}

\begin{theorem}[Two-sided relative curvature and smoothness]
\label{thm:ensemble-relative-smoothness}
Suppose Assumptions~\ref{ass:general-first-variation} and
\ref{ass:mirror-map} hold. For all $u,v\in\mathcal U[0,T]$,
\begin{equation}
\label{eq:ensemble-relative-smoothness}
\begin{aligned}
dJ^\tau[u](v-u)+(\tau-\kappa_{\mathrm{stab}})\mathcal D_h[v\mid u]
&\le J^\tau[v]-J^\tau[u]\\
&\le dJ^\tau[u](v-u)+(\tau+\kappa_{\mathrm{stab}})\mathcal D_h[v\mid u].
\end{aligned}
\end{equation}
The explicit relative-smoothness constant is
\begin{equation}
\label{eq:Lrel-general}
L_{\mathrm{rel}}:=\tau+\kappa_{\mathrm{stab}}.
\end{equation}
The upper bound is equivalently
\begin{align}
\label{eq:ensemble-relative-smoothness-H}
J^\tau[v]-J^\tau[u]\le{}&-\int_0^T\left\langle
\nabla_\theta H_t^\tau(x^u(t,\cdot),p^u(t,\cdot),u(t)),
v(t)-u(t)\right\rangle\,\dd t +L_{\mathrm{rel}}\mathcal D_h[v\mid u].
\end{align}
Define
\begin{equation}
\label{eq:mu-stab-general}
\mu_{\mathrm{stab}}:=\tau-\kappa_{\mathrm{stab}}.
\end{equation}
If $\tau=\kappa_{\mathrm{stab}}$, then $J^\tau$ is relatively convex;
if $\tau>\kappa_{\mathrm{stab}}$, it is relatively strongly convex
with modulus $\mu_{\mathrm{stab}}$.
\end{theorem}

\begin{proof}
Fix $u,v$ and set $\delta u=v-u$, $u^s=u+s\delta u$ for
$0\le s\le1$.
We first justify integrating the first variation along the
segment $s\mapsto u^s$. Set $\varphi(s):=J^\tau[u^s]$ and
\[
\psi(s):=-\int_0^T
\langle \nabla_\theta H_t^\tau(x^{u^s}(t,\cdot),p^{u^s}(t,\cdot),u^s(t)),\delta u(t)\rangle\,\dd t,
\qquad 0\le s\le1.
\]
For $\varepsilon>0$ we have
$u^s+\varepsilon\delta u=u^s+\frac{\varepsilon}{1-s}(v-u^s)$ when
$s\in[0,1)$, and
$u^s-\varepsilon\delta u=u^s+\frac{\varepsilon}{s}(u-u^s)$ when
$s\in(0,1]$. Applying Lemma~\ref{lem:first-variation-general} at the
base point $u^s$ in the feasible directions $v-u^s=(1-s)\delta u$ and
$u-u^s=-s\delta u$ therefore shows that $\varphi$ has right derivative
$\psi(s)$ at every $s\in[0,1)$ and left derivative $\psi(s)$ at every
$s\in(0,1]$. Hence $\varphi$ is continuous on $[0,1]$ and
differentiable on $(0,1)$ with $\varphi'=\psi$. Moreover, with
$F_\theta:=\sup|\nabla_\theta f|$ over the region used in
\eqref{eq:ensemble-convergence-bounds},
formula~\eqref{eq:general-Hamiltonian-u-derivative} gives
\[
|\nabla_\theta H_t^\tau(\xi,\pi,\theta)|
\le M_Q:=B_\theta M_P+F_\theta+\tau\max_{\vartheta\in U}|\nabla h(\vartheta)|,
\qquad \xi\in\mathcal X,\ \pi\in\mathcal P,\ \theta\in U,
\]
so $|\psi(s)|\le M_Q\sqrt T\,\|\delta u\|_{L^2}$. By the mean-value
theorem, $\varphi$ is Lipschitz on $[0,1]$, hence absolutely
continuous; $\psi$ is measurable as its a.e.\ derivative; and
$\varphi(1)-\varphi(0)=\int_0^1\psi(s)\,\dd s$. In other words,
\begin{equation}
\label{eq:J-segment-integral}
J^\tau[v]-J^\tau[u]=-\int_0^1\int_0^T
\langle \nabla_\theta H_t^\tau(x^{u^s}(t,\cdot),p^{u^s}(t,\cdot),u^s(t)),\delta u(t)\rangle\,\dd t\,\dd s.
\end{equation}
Substituting the first variation at $u$ and rearranging gives
\begin{align*}
J^\tau[v]-J^\tau[u]-dJ^\tau[u](v-u) = &-\int_0^1\int_0^T
\langle \nabla_\theta H_t^\tau(x^{u^s}(t,\cdot),p^{u^s}(t,\cdot),u^s(t))\\
&\qquad-\nabla_\theta H_t^\tau(x^u(t,\cdot),p^u(t,\cdot),u(t)),\delta u(t)\rangle\,\dd t\,\dd s.
\end{align*}
Since $H_t^\tau=H_t^0-\tau h$, the right-hand side splits into $I_2 -I_1$, where
\[
I_1 := \int_0^1\int_0^T
\langle \nabla_\theta H_t^0(x^{u^s}(t,\cdot),p^{u^s}(t,\cdot),u^s(t))-\nabla_\theta H_t^0(x^u(t,\cdot),p^u(t,\cdot),u(t)),\delta u(t)\rangle\,\dd t\,\dd s,
\]
and
\[
I_2 := \tau \int_0^1\int_0^T
\langle \nabla h(u^s(t)) - \nabla h(u(t)), \delta u(t) \rangle \, \dd t \, \dd s.
\]
Both are iterated integrals, with the time integral inside. The
inner integral of $I_2$ is continuous in $s$ by dominated convergence.
The inner integral of $I_1$ equals the inner integral of $I_2$ minus
$\psi(s)-\psi(0)$, so it is measurable in $s$. Hence $I_1$ and $I_2$
are well defined.
By the fundamental theorem of calculus applied to the convex function $h$, we have
\[
I_2 =\tau\mathcal D_h[v\mid u].
\]
Thus the remainder identity is
\begin{equation}
\label{eq:relative-smoothness-remainder}
J^\tau[v]-J^\tau[u]-dJ^\tau[u](v-u)
=-I_1+\tau\mathcal D_h[v\mid u].
\end{equation}
The two stability lemmas and Cauchy--Schwarz give
\begin{align*}
|I_1|&\le C_H\int_0^1
\bigl(\|x^{u^s}-x^u\|_{L^2}+\|p^{u^s}-p^u\|_{L^2}
+\|u^s-u\|_{L^2}\bigr)\|\delta u\|_{L^2}\,\dd s\\
&\le\frac{C_H(C_x+C_p+1)}2\|\delta u\|_{L^2}^2\\
&\le\kappa_{\mathrm{stab}}\mathcal D_h[v\mid u],
\end{align*}
where the last step uses
$\mathcal D_h[v\mid u]\ge(\sigma_h/2)\|\delta u\|_{L^2}^2$.
Hence $-\kappa_{\mathrm{stab}}\mathcal D_h[v\mid u]\le-I_1
\le\kappa_{\mathrm{stab}}\mathcal D_h[v\mid u]$.
Substitution into \eqref{eq:relative-smoothness-remainder} proves both
sides of \eqref{eq:ensemble-relative-smoothness}.
\end{proof}

\end{document}
